\documentclass{article}
\usepackage{graphicx} % Required for inserting images


% ---------- packages ----------
\usepackage[margin=1in]{geometry}
\usepackage{amsmath, amssymb, amsthm, bm, mathtools}
\usepackage{enumitem}
\usepackage{hyperref}
\usepackage{booktabs}
\usepackage{microtype}

% ---------- macros ----------
\newcommand{\R}{\mathbb{R}}
\newcommand{\Tkn}{\mathcal{T}}
\newcommand{\Flow}{\Phi}
\newcommand{\Attn}{\mathrm{attn}}
\newcommand{\MLP}{\mathrm{mlp}}
\newcommand{\Fric}{\mathrm{fric}}
\newcommand{\Lie}{\mathrm{Lie}}
\newcommand{\Strang}{\mathrm{Strang}}
\newcommand{\tol}{\mathrm{tol}}
\newcommand{\NFE}{\mathrm{NFE}}
\newcommand{\diag}{\mathrm{diag}}
\newcommand{\half}{\tfrac{1}{2}}

\title{\textbf{Structure-Aware Split-Flow ODE Solver}}
\author{jeon.isavelle@gmail.com}
\date{\today}

\begin{document}
\maketitle

\begin{abstract}
We propose a structure-aware split-flow ODE solver for ViTs. Instead of applying a black-box ODE solver to the full transformer dynamics, we decompose the vector field into attention-induced mixing, token-wise MLP reaction, and friction/retention flows. The resulting solver is implemented through operator splitting, optionally with adaptive step-size control based on split-pair error estimation. This design aims to improve stability, and to compute allocation during inference.
\end{abstract}

\section{Introduction}

Vision Transformers (ViTs) are effective for modeling global context; however, in 3D volumes, the number of tokens grows rapidly, leading to a substantial computational burden. This issue becomes particularly important in medical image segmentation, where volumetric inputs require both fine spatial reasoning and long-range contextual modeling. As a result, directly applying standard ViT blocks to 3D data often yields high computational cost and limited flexibility in allocating computation according to the difficulty of the input.

To address this issue, this report considers a structure-aware adaptive computation framework based on an ODE interpretation of transformer blocks. Rather than relying solely on a generic black-box continuous-depth formulation, we design a split-flow solver that explicitly reflects the internal structure of the transformer. In particular, the proposed approach decomposes the dynamics into attention, MLP, and retention/friction flows, and connects this decomposition to an operator-splitting-based numerical scheme. Within this formulation, adaptive step-size control serves as a mechanism for adjusting computational effort according to the complexity of the latent dynamics, providing a principled link between continuous-time modeling and efficient inference in 3D segmentation.

\paragraph{Main idea}
Let the token state be
\begin{equation}
    u(t) \in \R^{T \times d},
\end{equation}
where $T$ is the number of tokens and $d$ is the embedding dimension. Rather than defining a single undifferentiated vector field $u' = f(u,t)$, we design a structured decomposition that reflects the internal organization of the transformer block.

\paragraph{Contributions}
\begin{enumerate}[leftmargin=1.5em]
    \item A three-way decomposition of the transformer dynamics into attention, MLP, and friction flows.
    \item A Lie--Trotter / Strang split solver tailored to this decomposition (A symmetric second-order split solver tailored to this decomposition three operators).
    \item An optional adaptive controller based on split-pair local error estimation.
    \item A practical stability interpretation via bounded gates and contraction-like retention.
\end{enumerate}

\section{Method}

\subsection{Decomposition of the Vector Field}
We decompose the dynamics as
\begin{equation}
    \frac{du}{dt} = f_{\Attn}(u,t) + f_{\MLP}(u,t) + f_{\Fric}(u).
\end{equation}
Here denotes:
\begin{itemize}[leftmargin=1.5em]
    \item $f_{\Attn}$ the attention-induced token mixing dynamics,
    \item $f_{\MLP}$ the token-wise nonlinear reaction dynamics,
    \item $f_{\Fric}$ the retention/friction term that stabilizes the evolution.
\end{itemize}

We adopt the retention--forcing form
\begin{equation}
    u' + P(u)u = Q(u),
\end{equation}
which can be rewritten as
\begin{equation}
    u' = -P(u)u + Q(u).
\end{equation}
We further split the forcing term into
\begin{equation}
    Q(u) = Q_{\Attn}(u) + Q_{\MLP}(u),
\end{equation}
so that
\begin{equation}
    f_{\Fric}(u) = -P(u)u,
    \qquad
    f_{\Attn}(u,t) = Q_{\Attn}(u,t),
    \qquad
    f_{\MLP}(u,t) = Q_{\MLP}(u,t).
\end{equation}

\paragraph{Design note.}
To limit drift and explosion, we use bounded gates for the coefficients in $P$ and the forcing terms.

\subsection{Sub-Flow Maps}
Each sub-vector field induces a flow map over step size $h$:
\begin{equation}
    \phi_{\Attn}^{h},
    \qquad
    \phi_{\MLP}^{h},
    \qquad
    \phi_{\Fric}^{h}.
\end{equation}
The friction flow can be treated in a particularly simple manner. If $P(u)$ is frozen over a substep, then
\begin{equation}
    u' = -\widehat{P}u
\end{equation}
has the closed-form update
\begin{equation}
    u_{n+1} = \exp(-h\widehat{P})u_n.
\end{equation}
For diagonal or channel-wise $\widehat{P}$, this update is cheap and explicitly damping.

\subsection{Split Solver}
Given that the proposed dynamics are decomposed into three components, the research focuses on implementing a symmetric second-order splitting adapted to the three-way decomposition. Let
\begin{equation}
    A := f_{\Fric},
    \qquad
    B := f_{\Attn},
    \qquad
    C := f_{\MLP}.
\end{equation}

A first-order Lie--Trotter approximation is given by
\begin{equation}
    \Flow_h^{\Lie}
    \approx
    \phi_{\MLP}^{h}
    \circ
    \phi_{\Attn}^{h}
    \circ
    \phi_{\Fric}^{h}.
\end{equation}
A symmetric second-order Strang-style solver can be written as
\begin{equation}
    \Flow_h^{\Strang}
    \approx
    \phi_{\Fric}^{h/2}
    \circ
    \phi_{\Attn}^{h/2}
    \circ
    \phi_{\MLP}^{h}
    \circ
    \phi_{\Attn}^{h/2}
    \circ
    \phi_{\Fric}^{h/2}.
\end{equation}

Equivalently, in exponential notation,
\begin{equation}
    \Flow_h^{\Strang}
    \approx
    e^{\frac{h}{2}A}
    e^{\frac{h}{2}B}
    e^{hC}
    e^{\frac{h}{2}B}
    e^{\frac{h}{2}A}.
\end{equation}

\paragraph{Why splitting?}
The split solver makes the role of each mechanism explicit and allows distinct numerical treatment for mixing, reaction, and contraction. In practice, this can improve the stability and accuracy trade-off relative to a single black-box solver under the same computational budget.

\subsection{Adaptive Step-Size Control}

To support adaptive computation, we estimate local error through a split-pair controller. Let
\begin{equation}
    u_{n+1}^{(1)} = \Flow_h^{\Lie}(u_n),
    \qquad
    u_{n+1}^{(2)} = \Flow_h^{\Strang}(u_n).
\end{equation}
We define the normalized local error surrogate as
\begin{equation}
    e_n
    =
    \frac{\lVert u_{n+1}^{(2)} - u_{n+1}^{(1)} \rVert}
    {\mathrm{atol} + \mathrm{rtol}\,\lVert u_{n+1}^{(2)} \rVert}.
\end{equation}
If $e_n \le 1$, the step is accepted; otherwise the step is rejected and recomputed with a smaller $h$.

\section{Implementation}

\subsection{Block-Level Update}
One ODE block applies the following sequence:
\begin{enumerate}[leftmargin=1.5em]
    \item compute or freeze friction coefficients $P(u_n)$,
    \item apply a half or full friction update,
    \item apply attention subflow on the current token state,
    \item apply MLP subflow token-wise,
    \item recombine according to Lie or Strang ordering.
\end{enumerate}


\subsection{Practical Options}
\begin{itemize}[leftmargin=1.5em]
    \item \textbf{Friction parameterization:} scalar, channel-wise diagonal, or token-wise diagonal.
    \item \textbf{Attention subsolver:} Euler, midpoint, or short RK2 on a frozen attention operator.
    \item \textbf{MLP subsolver:} explicit residual step or token-wise micro-step update.
    \item \textbf{Adaptive controller:} enabled during inference only, disabled during training if needed for stability.
\end{itemize}

\section{Expected Advantages}
\begin{itemize}[leftmargin=1.5em]
    \item \textbf{Interpretability:} separates token mixing, token-local reaction, and retention.
    \item \textbf{Stability control:} friction acts as an explicit contraction-like mechanism.
    \item \textbf{Compute allocation:} adaptive step control allows variable computational effort.
    \item \textbf{Modularity:} each subflow can be assigned a solver appropriate to its structure.
\end{itemize}

\section{Limitations and Caveats}
\begin{itemize}[leftmargin=1.5em]
    \item Attention and MLP flows are generally non-commutative, so splitting error may be nontrivial.
    \item If all tokens are evolved as one global state, adaptivity remains mostly global rather than strictly region-wise.
    \item Excessively weak friction may fail to stabilize the dynamics; excessively strong friction may oversmooth features.
    \item The method improves numerical control, but does not by itself constitute a full convergence proof.
\end{itemize}

\section{Conclusion}
TODO

\section{Future Work}
The further research would include:

\begin{itemize}[leftmargin=1.5em]
    \item complete appendix sector,
    \item exploring alternative symmetric second-order orderings for the three-operator decomposition,
    \item comparing fixed-step and adaptive split solvers empirically.
\end{itemize}

% ---------- optional appendix ----------
\appendix
\section{Notation Checklist}
\begin{itemize}[leftmargin=1.5em]
    \item $T$: number of tokens
    \item $d$: embedding dimension
    \item $u(t) \in \R^{T \times d}$: token state
    \item $h$: solver step size
    \item $P(u)$: friction/retention coefficient map
    \item $Q_{\Attn}, Q_{\MLP}$: attention and MLP forcing terms
\end{itemize}

\end{document}
